\subsection{RTL Controller and Clock-Gating Organization}
\label{sec:rtl-clock-gating}

We implement the causal predictor as a synthesizable SystemVerilog controller
with 20-bit horizon and memory-ticket fields.  The controller stores one
descriptor for each of the two resident tile slots.  A descriptor contains
tile presence, the first future SA opcode and its prefix distance, remaining
tile work, and the allocation, validity, pending-memory, tail-ticket, and
pending-response state of at most two source operands.  The combinational
bound logic computes the prefix, request-FIFO, response-drain, visible-work,
and future-tile bounds described in
Secs.~\ref{sec:causal-predictor} and
\ref{sec:causal-bound-components}.

The implementation separates this state into three clock domains.  First, a
descriptor domain captures resident-tile metadata only on an explicit write
to either tile slot.  Second, an evaluation domain implements the three-stage
certificate pipeline.  Third, a small control domain remains on the root
clock so that it can commit an entry, decrement the frozen sleep countdown,
and apply the immediate-wake conditions.  Keeping the control domain
available avoids relying on a clock that the controller itself may stop.

The evaluation pipeline is activated only by a source-grounded sampling
event.  A new sample is requested when the complete entry-eligibility
predicate changes from false to true, after a semantic descriptor update, or
after the countdown domain exits a sleep interval.  Once a sample enters the
pipeline, registered valid bits hold its clock enable active through bound
calculation and entry commit.  Four enabled edges are sufficient to capture
stage~1, propagate stages~2 and~3, and commit the result.  Per-cycle dequeue
progress and pending-response-count reductions do not independently restart
the evaluation pipeline because they can only shorten the previously observed
lower bound.  Their current values are sampled when a fresh certificate is
constructed.

The wrapper instantiates two ASAP7 integrated clock-gating cells and two local
clock buffers.  One gate--buffer pair drives the certificate-evaluation
domain; the other drives the descriptor-register domain.  The implementation
uses \texttt{ICGx2\_ASAP7\_75t\_R} clock gates and
\texttt{BUFx24\_ASAP7\_75t\_R} buffers.  The synthesized netlist retains
exactly two integrated clock gates and two clock buffers.  The small event
detector and countdown/wake logic remain on the ungated root clock.

Clock gating does not change the certificate acceptance rule.  The wrapped
controller receives the same architectural descriptor and queue state as the
continuously evaluated reference controller.  Descriptor writes keep the
descriptor clock enabled and leave a pending event that initiates evaluation
only after the updated descriptor has been registered.  A certificate-stage
valid bit keeps the evaluation clock enabled until the in-flight sample has
been consumed.  The resulting \texttt{sleep\_request} is therefore produced
from a complete registered certificate rather than from a partially updated
descriptor.

The RTL boundary includes the descriptor registers, causal-bound arithmetic,
certificate pipeline, clock-gating wrapper, event detector, and countdown
control.  The current ONNXim trace generator constructs resident descriptors
by scanning software instruction deques.  A hardware descriptor producer
that replaces this software scan is outside the synthesized controller
boundary; consequently, its area and power are not included in the reported
controller implementation.  Likewise, the controller emits the requested
SA power state, but isolation cells, a physical SA power switch, and their
power-distribution network are evaluated separately rather than instantiated
inside this RTL block.

\subsection{Replay Equivalence and Safety Validation}
\label{sec:rtl-replay-validation}

We validate the clock-gated RTL using cycle-indexed architectural traces from
the same workload executions used in the evaluation.  Each trace begins with
a complete cycle-zero record and then contains state-change records.  The
testbench holds each record until the next recorded cycle, thereby
reconstructing the controller inputs for every intervening cycle.  It rejects
an empty trace, a malformed record, or a non-monotonic cycle sequence.

Each core trace is replayed twice: once through the continuously evaluated
reference controller and once through the clock-gated wrapper.  For both
runs, the validation flow records the trace marker, every sleep episode, and
the final replay summary.  The two emitted record streams are compared
byte-for-byte.  Any difference in sleep entry, sleep exit, episode duration,
transition count, or summary field causes the comparison to fail.  A workload
is accepted only after all four core traces pass this comparison.

The testbench additionally increments a safety-violation counter on every
cycle satisfying
\begin{equation}
  \texttt{sleep\_request} \land \texttt{sa\_active}.
  \label{eq:replay-safety-violation}
\end{equation}
It terminates with an error when this counter is nonzero.  This check directly
tests that a requested non-retentive OFF interval never overlaps SA activity
observed in the architectural trace.

Table~\ref{tab:rtl-replay-validation} summarizes the retained replay records
for the 16 workload configurations in Table~\ref{tab:component-characterization}.
The totals are sums of the four per-core summaries for every workload.

\begin{table}[t]
  \centering
  \caption{Aggregate RTL replay-validation results.}
  \label{tab:rtl-replay-validation}
  \begin{tabular}{lr}
    \toprule
    Metric & Measured value \\
    \midrule
    Workload configurations & 16 \\
    Core traces & 64 \\
    Replayed core-cycles & 898,576,188 \\
    Requested sleep cycles & 191,126,583 \\
    Sleep-state transitions & 1,674,074 \\
    Completed sleep episodes & 837,037 \\
    Sleep--SA-activity violations & 0 \\
    \bottomrule
  \end{tabular}
\end{table}

All 64 retained core replays produced identical reference and clock-gated
sleep-record streams, and all reported zero violations of
Eq.~\eqref{eq:replay-safety-violation}.  These results establish equivalence
of the two controller implementations for the evaluated traces and show that
their requested OFF intervals do not overlap trace-observed SA activity.

The scope of this result is intentionally limited.  Byte-identical replay is
not a formal proof over all possible input sequences, and the implemented
safety counter checks sleep--activity overlap rather than every internal
invariant.  Moreover, power-state requests are evaluated in shadow replay and
do not alter ONNXim scheduling or insert a physical wake-up delay.  Therefore,
the replay result does not by itself establish zero performance penalty for a
fully integrated power-gated NPU.  Such a claim would require feeding the
measured transition latency back into architectural execution and verifying
the resulting timing behavior.
